\section{核心模块实现}

\subsection{配置管理模块 (config.py)}

配置管理模块负责加载和管理系统的配置参数，确保所有模块使用统一的配置。

\subsubsection{功能实现}
- 从 `configs/config.yaml` 文件加载配置
- 提供全局单例配置对象 `GLOBAL_CONFIG`，供其他模块直接导入使用
- 支持配置文件路径的自定义

\subsubsection{关键代码}

\begin{lstlisting}[language=Python]
def load_config(config_path: str | Path = DEFAULT_CONFIG_PATH) -> Dict[str, Any]:
    if not os.path.exists(config_path):
        raise FileNotFoundError(f"Config file {config_path} not found.")
    try:
        with open(config_path, "r", encoding="utf-8") as f:
            config = yaml.safe_load(f)
            return config
    except yaml.YAMLError as e:
        raise ValueError(f"YAML error in {config_path}: {e}")

GLOBAL_CONFIG = load_config()
\end{lstlisting}

\subsection{数据接入模块 (ingest)}

数据接入模块负责 PDF 研报的解析、分块和向量化，为后续的检索和生成提供基础数据。

\subsubsection{PDF 解析 (pdf_parser.py)}

\paragraph{功能实现}
- 使用 PyMuPDFReader 解析 PDF 研报，提取文本内容
- 确保每个文档的元数据中包含 `source` 和 `page_label`
- 规范化元数据格式，确保符合系统的数据契约

\paragraph{关键代码}

\begin{lstlisting}[language=Python]
def load_financial_pdfs(data_dir: str | Path = None) -> List[Document]:
    parser = PyMuPDFReader()
    for pdf_path in pdf_files:
        file_name = pdf_path.name
        try:
            docs_from_file = parser.load_data(file_path=str(pdf_path))
            for i, doc in enumerate(docs_from_file):
                raw_metadata = doc.metadata or {}
                extracted_page = raw_metadata.get("page", raw_metadata.get("source_page", str(i + 1)))
                doc.metadata = {
                    "source": file_name,
                    "page_label": str(extracted_page),
                    "total_pages": raw_metadata.get("total_pages", "unknown"),
                }
                doc.excluded_embed_metadata_keys = ["page_label", "total_pages", "source"]
                doc.excluded_llm_metadata_keys = ["total_pages"]
                all_documents.append(doc)
        except Exception as e:
            print(f"Failed: {file_name} - {str(e)}")
    return all_documents
\end{lstlisting}

\subsubsection{分块策略 (chunker.py)}

\paragraph{功能实现}
- 实现基线分块策略：使用固定大小的分块，块大小为 256 tokens
- 实现语义分块策略：基于 BGE-M3 的语义分块，通过计算句间余弦相似度进行切分

\subsubsection{索引构建 (indexer.py)}

\paragraph{功能实现}
- 使用 BGE-M3 模型对分块后的文本进行向量化
- 将向量存储到本地 ChromaDB 中
- 支持 RAPTOR 树状摘要的构建

\subsection{检索模块 (retrieval)}

检索模块负责根据用户查询从向量库中检索相关信息，为生成模块提供上下文。

\subsubsection{检索策略 (retriever.py)}

\paragraph{功能实现}
- 实现单路稠密向量检索（Phase 1）
- 实现多路混合检索（Phase 2）：集成 BM25 检索和向量检索，通过 RRF 融合结果
- 实现精排：使用 bge-reranker-base 对检索结果进行精排
- 实现全局单例缓存，提高性能

\paragraph{关键代码}

\begin{lstlisting}[language=Python]
def get_hybrid_retriever(index: VectorStoreIndex) -> BaseRetriever:
    global _CACHED_BM25_RETRIEVER
    vector_top_k = GLOBAL_CONFIG["retrieval"].get("vector_top_k", 20)
    bm25_top_k = GLOBAL_CONFIG["retrieval"].get("bm25_top_k", 20)
    vector_retriever = index.as_retriever(similarity_top_k=vector_top_k)
    if _CACHED_BM25_RETRIEVER is None:
        print("Building BM25 retriever...")
        nodes_for_bm25 = _extract_nodes_from_vector_store(index)
        if not nodes_for_bm25:
            return vector_retriever
        _CACHED_BM25_RETRIEVER = BM25Retriever.from_defaults(
            nodes=nodes_for_bm25, similarity_top_k=bm25_top_k
        )
    fusion_retriever = QueryFusionRetriever(
        retrievers=[vector_retriever, _CACHED_BM25_RETRIEVER],
        similarity_top_k=GLOBAL_CONFIG["retrieval"].get("vector_top_k", 20),
        num_queries=1,
        mode="reciprocal_rerank",
        llm=Ollama(
            model=GLOBAL_CONFIG["llm"]["weak_model"],
            base_url=GLOBAL_CONFIG["llm"]["ollama_base_url"],
            temperature=0.0,
            request_timeout=300.0,
        ),
    )
    return fusion_retriever
\end{lstlisting}

\subsubsection{精排模块 (reranker.py)}

\paragraph{功能实现}
- 使用 bge-reranker-base 对检索结果进行精排
- 实现全局单例缓存，提高性能

\subsection{生成模块 (generation)}

生成模块负责构建对话引擎和生成回答，为用户提供智能的金融研报分析。

\subsubsection{对话引擎 (pipeline.py)}

\paragraph{功能实现}
- 使用 CondensePlusContextChatEngine 作为核心对话引擎
- 集成 ChatMemoryBuffer 实现短期记忆管理
- 支持多轮对话，理解上下文指代

\subsubsection{LLM 后端 (llm_backend.py)}

\paragraph{功能实现}
- 封装 Ollama 本地部署的大语言模型
- 支持不同模型的切换

\subsubsection{提示词设计 (prompt.py)}

\paragraph{功能实现}
- 设计严格的反幻觉和强制引用指令
- 确保生成的回答包含引用信息

\subsection{评估模块 (evaluation)}

评估模块负责对系统的性能进行评估，确保系统输出的质量。

\subsubsection{RAGAS 评估 (ragas_eval.py)}

\paragraph{功能实现}
- 使用 RAGAS 评估框架对系统回答进行评估
- 将本地运行的 Qwen2.5-72B 注册为 RAGAS 的 judge_model
- 计算 Context Precision、Faithfulness 和 Answer Relevancy 指标

\subsubsection{引用审计 (citation_audit.py)}

\paragraph{功能实现}
- 审计生成回答中的引用信息
- 确保引用的页码和源文件存在

\subsection{交互模块 (ui)}

交互模块负责构建用户界面，提供友好的用户体验。

\subsubsection{前端界面 (app.py)}

\paragraph{功能实现}
- 使用 Gradio 构建前端界面
- 实现多文档管理、对话式问答、引用跳转和产物生成功能
- 支持引用点击跳转到 PDF 对应页码
